\pdfbookmark[0]{Notices. Monge: Géométrie descriptive. 1813.}{anzeigen-monge-1813}
\begin{center}
\rule{.85\textwidth}{.3pt}\par
Göttingische gelehrte Anzeigen.\quad 31 July 1813.\par
\rule{.85\textwidth}{.3pt}
\end{center}
\hspace*{\parindent}\textit{Géométrie descriptive par} \textsc{Gaspard Monge}, \textit{de l’institut des sciences etc. Nouvelle édition. Avec un supplément par} \textsc{M. Hachette}, \textit{instituteur à l’école impériale polytechnique etc.} Paris, published by \textsc{J. Klostermann} the younger. 162 and 118 pages, in quarto.

Geometry studies spatial relations. It divides into two broad branches depending on whether space is considered in only two dimensions, in the plane, or in all three dimensions together. It is easy to see that the second branch must by its nature cover much more ground and offer a far greater variety of questions and inquiries than the first. If even in our elementary geometry plane geometry takes up more room than solid geometry, this is simply because the latter has been worked out and developed much less fully. Indeed, in more recent times there has been a preference for treating inquiries of the latter sort in particular by analysis, thus in a sense taking them away from geometry, which relies entirely on direct intuition. The advantages of the analytical over the geometrical treatment are undeniable. Its brevity, simplicity, uniform procedure and especially its generality tend to stand out all the more clearly as the inquiries become harder and more intricate. Even so, it remains very important to keep developing the geometrical method as well. In some particular %
